% Part: first-order-logic
% Chapter: tableaux
% Section: rules-and-proofs

\documentclass[../../../include/open-logic-section]{subfiles}

\begin{document}

\iftag{FOL}
      {\olfileid{fol}{tab}{rul}}
      {\olfileid{pl}{tab}{rul}}

\olsection{Aturan lan \usetoken{P}{tableau}}

!!^a{tableau} yaiku pamriksa sistematis tumrap kabeh cara kang bisa
ndadekake !!a{sentence} bener utawa salah ing !!a{structure}. Unsur
pambangun tableau yaiku !!{signed formula}s: !!{sentence}s bebarengan
karo ``tandha'' nilai kabeneran, yaiku salah siji saka $\True$
utawa~$\False$. !!{formula}s iki disusun dadi wit (kang tuwuh
mudhun).

\begin{defn}
  \emph{!!{signed formula}} yaiku pasangan kang dumadi saka nilai
  kabeneran lan !!a{sentence}, yaiku salah siji saka:
  \[
  \sFmla{\True}{!A} \text{ utawa } \sFmla{\False}{!A}.
  \]
\end{defn}

Miturut intuisi, $\sFmla{\True}{!A}$ bisa diwaca minangka ``$!A$
bisa wae bener'' lan $\sFmla{\False}{!A}$ minangka ``$!A$ bisa wae
salah'' (ing sawijining !!{structure}).

Saben !!{signed formula} ing wit iku salah siji saka \emph{asumsi}
(kang didhaptar ing pucuk wit), utawa diasilake saka
!!a{signed formula} ing ndhuwure liwat salah siji saka sawetara aturan inferensi.
Ana rong aturan kanggo saben !!{main operator} kang bisa diduweni
!!{formula} sadurunge, siji kanggo tandha~$\True$ lan siji kanggo
tandha~$\False$. Sawetara aturan ngidini wit nyabang, dene liyane mung
nambahake !!{signed formula}s menyang sawijining cabang. Sawijining
aturan bisa (lan kerep kudu) ditrapake dudu marang !!{signed formula}
kang langsung ana sadurunge, nanging marang sembarang
!!{signed formula} ing cabang saka oyod nganti panggonan aturan iku ditrapake.

Sawijining cabang \emph{katutup} yen ngemot $\sFmla{\True}{!A}$ lan
$\sFmla{\False}{!A}$ bebarengan. !!{tableau} katutup yen saben
cabange katutup. Miturut interpretasi intuisi mau, saben cabang
nggambarake sawijining kamungkinan bebarengan, nanging
$\sFmla{\True}{!A}$ lan $\sFmla{\False}{!A}$ ora bisa kelakon
bebarengan. Kanthi tembung liya, yen cabang katutup, kamungkinan kang
digambarake wis kasisih. Mligine, iki ateges !!{tableau} katutup
nyingkirake kabeh kamungkinan kanggo bebarengan ndadekake saben asumsi
awujud $\sFmla{\True}{!A}$ bener lan saben asumsi awujud
$\sFmla{\False}{!A}$ salah.

!!{tableau} katutup \emph{kanggo $!A$} yaiku !!{tableau} katutup
kanthi oyod~$\sFmla{\False}{!A}$. Yen ana !!{tableau} katutup kaya
mangkono, kabeh kamungkinan yen~$!A$ salah wis kasisih; yaiku, $!A$
kudu bener ing saben !!{structure}.

\end{document}
